% English abstract
\chapter*{Abstract}
\addcontentsline{toc}{chapter}{Abstract}
\begin{latin}
Numerical solution of nonlinear partial differential equations is a
fundamental problem in scientific computing. Classical numerical methods,
while accurate and widely used, depend on computational meshes and struggle
in high dimensions; on the other hand, purely data-driven deep learning
models require large amounts of training data and offer no guarantee that
their outputs respect the laws of physics. In recent years, the framework
of physics-informed neural networks has been introduced as a way to combine
physical knowledge (the governing equation of the system) with machine
learning. The aim of this thesis is to study, implement, and evaluate this
framework for solving forward and inverse problems involving nonlinear
partial differential equations.

In this work, we first review the required theoretical background, including
partial differential equations, classical numerical methods, neural networks,
automatic differentiation, and the background of physics-aware machine
learning. We then describe the physics-informed framework, including the
definition of the residual network, the combined data-physics loss function,
and the continuous- and discrete-time models. In the practical part, the
continuous-time model is implemented for the one-dimensional Burgers
equation in Python using the JAX library, and an independent reference
solver based on the Fourier spectral method and fourth-order time stepping
is developed to generate the exact solution and verified against the
Cole-Hopf semi-analytical solution.

In the forward problem, the solution over the whole spatio-temporal domain
is reconstructed using only 200 initial and boundary data points, achieving
a relative error of \FwdRelErr\ in the $L_2$ norm; the error is mostly
concentrated in the thin shock layer. In the inverse problem, the equation
coefficients are identified from 5000 scattered observations: the advection
coefficient with an error of \LamIErrClean\ and the diffusion coefficient
with an error of \LamIIErrClean\ in the noise-free case, and \LamIErrNoisy\
and \LamIIErrNoisy\ respectively under one percent noise. The results show
that this framework achieves an acceptable solution with little data, is
reasonably robust to noise, and that its main bottleneck is resolving the
steep solution layers and the inherent sensitivity of identifying small
coefficients.

\textbf{Keywords:} Physics-informed neural network, Nonlinear partial
differential equation, Forward problem, Inverse problem, Burgers equation,
Automatic differentiation
\end{latin}
